\documentclass[11pt]{article}
\usepackage[margin=1in]{geometry}
\usepackage{amsmath,amssymb,amsthm,mathtools}
\usepackage{graphicx}
\usepackage{xcolor}
\usepackage{microtype}
\usepackage[colorlinks=true,linkcolor=blue!55!black,citecolor=blue!55!black,urlcolor=blue!55!black]{hyperref}

\newtheorem{theorem}{Theorem}
\newtheorem{lemma}[theorem]{Lemma}
\newtheorem{remark}[theorem]{Remark}
\newcommand{\R}{\mathbb R}
\newcommand{\T}{\mathbb T}
\newcommand{\Z}{\mathbb Z}
\newcommand{\cF}{\mathcal F}
\newcommand{\M}{\mathcal M}
\newcommand{\eps}{\varepsilon}
\newcommand{\wh}{\widehat}

\title{De Leeuw transference on homogeneous $\dot W^{1,1}$}
\author{Solution packet for arXiv:1306.1437v4}
\date{12 August 2026}

\begin{document}
\maketitle

\noindent\textbf{Status.}
Candidate full solution in the exact space of the source paper, likely valid,
awaiting human review.

\section{Source question and result}

Kazaniecki and Wojciechowski prove that every Fourier multiplier on
$\dot W^{1,1}(\R^d)$, $d\ge2$, has a bounded continuous representative
\cite{KW}.  Their argument needs estimates on the torus.  They observe that
ordinary $L^1$ has de Leeuw transference, but write that no equivalent was
known for the homogeneous Sobolev space and leave the general de Leeuw-type
question open.  The statement is on source PDF page~3:

\begin{center}
 \includegraphics[width=0.94\textwidth]{figures/open_problem_crop.png}
\end{center}

We prove both analogues of the classical de Leeuw theorems which the source
states immediately afterward: restriction to a lattice and approximation
from shrinking lattices.  In particular, we obtain an exact norm identity.

Put $\T^d=\R^d/(2\pi\Z^d)$ and use the seminorms
\[
 \|f\|_{\dot W_{\R}}=\sum_{j=1}^d\|\partial_jf\|_{L^1(\R^d)},
 \qquad
 \|P\|_{\dot W_{\T}}=\sum_{j=1}^d\|\partial_jP\|_{L^1(\T^d)}.
\]
Both spaces are understood modulo constants.  Multiplier norms are denoted
by $\|m\|_{\M(\dot W_{\R})}$ and
$\|\gamma\|_{\M(\dot W_{\T})}$.  The coefficient $\gamma_0$ never affects
the torus seminorm.

\subsection*{Proof Intuition}

For restriction, test the Euclidean multiplier on a trigonometric polynomial
inside a slowly varying band-limited envelope.  The apparent obstruction is
that a homogeneous $\dot W^{1,1}$ multiplier need not globally be convolution
by a measure.  All finitely many wave packets avoid frequency zero, however;
after a smooth localization there, one can invert a coordinate derivative
and turn the multiplier into an ordinary $L^1$ measure multiplier.  Translation
continuity then identifies the torus action with no loss in the limiting
constant.  In the reverse direction, Poisson summation turns inverse-Fourier
Riemann sums into periodizations of a band-limited Schwartz function.  The
input norms converge on expanding period cubes, while the output gradients
converge locally to the Euclidean multiplier output.  The two limits give the
reverse inequality and hence exact norm equality.

\begin{theorem}[Homogeneous de Leeuw theorem]\label{thm:main}
Let $d\ge2$ and let $m:\R^d\to\mathbb C$ be continuous.  For $\eps>0$ set
\[
 \gamma_k^{\eps}=m(\eps k),\qquad k\in\Z^d.
\]
Then the following statements hold.
\begin{enumerate}
\item If $m\in\M(\dot W_{\R})$, then
\[
 \gamma^\eps\in\M(\dot W_{\T}),\qquad
 \|\gamma^\eps\|_{\M(\dot W_{\T})}
 \le \|m\|_{\M(\dot W_{\R})}
 \quad(\eps>0).
\]
\item If $\eps_n\downarrow0$ and
$\sup_n\|\gamma^{\eps_n}\|_{\M(\dot W_{\T})}\le C$, then
\[
 m\in\M(\dot W_{\R}),\qquad
 \|m\|_{\M(\dot W_{\R})}\le C.
\]
\end{enumerate}
Consequently,
\begin{equation}\label{eq:norm-identity}
 \|m\|_{\M(\dot W_{\R})}
 =\sup_{\eps>0}\|(m(\eps k))_{k\in\Z^d}\|_{\M(\dot W_{\T})},
\end{equation}
with the convention that either side may be infinite.
\end{theorem}

Because the continuity hypothesis is automatic for Euclidean multipliers by
the main theorem of \cite{KW}, Theorem~\ref{thm:main} applies to every
multiplier in the precise setting of that paper.  It does not claim the same
for every homogeneous order and exponent; see the scope note at the end.

We use the Fourier convention
\[
 \wh f(\xi)=\int_{\R^d}f(x)e^{-ix\cdot\xi}\,dx,
 \qquad
 f(x)=(2\pi)^{-d}\int_{\R^d}\wh f(\xi)e^{ix\cdot\xi}\,d\xi.
\]

\section{Two localization lemmas}

The first observation repairs the obstacle in a direct cutoff argument:
although $m$ need not be a Fourier transform of a measure globally, every
smooth localization away from the origin is one.

\begin{lemma}[Local $L^1$ multiplier]\label{lem:local-L1}
Suppose $m\in\M(\dot W_{\R})$ and
$\chi\in C_c^\infty(\R^d\setminus\{0\})$.  Then $m\chi$ is an ordinary
$L^1(\R^d)$ multiplier.  Hence $m\chi=\wh\mu$ for a finite complex Borel
measure $\mu$.
\end{lemma}

\begin{proof}
Choose a finite smooth decomposition $\chi=\sum_r\chi_r$ such that, for
each $r$, some coordinate $\xi_{j_r}$ is bounded away from zero on
$\operatorname{supp}\chi_r$.  Given $g\in\mathcal S(\R^d)$, define
\[
 \wh f_r(\xi)=\frac{\chi_r(\xi)}{i\xi_{j_r}}\wh g(\xi).
\]
The multiplier identity gives
\[
 T_{m\chi_r}g=\partial_{j_r}T_mf_r.
\]
For every $\ell$, the symbol
$\xi_\ell\chi_r(\xi)/\xi_{j_r}$ is smooth and compactly supported.  Its
inverse Fourier transform is in $L^1$, so Young's inequality gives
\[
 \|f_r\|_{\dot W_{\R}}\le C_{\chi_r}\|g\|_1.
\]
It follows that
\[
 \|T_{m\chi_r}g\|_1
 \le \|T_mf_r\|_{\dot W_{\R}}
 \le \|m\|_{\M(\dot W_{\R})}C_{\chi_r}\|g\|_1.
\]
Summing in $r$ makes $m\chi$ an $L^1$ multiplier.  The classical
characterization of translation-invariant bounded operators on $L^1$
identifies their symbols with Fourier transforms of finite measures
\cite{Rudin}.
\end{proof}

For a finite set $F\subset\Z^d$, write
$T_{b|F}(\sum_{k\in F}c_ke^{ik\cdot x})
=\sum_{k\in F}b(k)c_ke^{ik\cdot x}$.

\begin{lemma}[Finite-frequency wave packets]\label{lem:packet}
Let $b=\wh\mu$ for a finite complex measure $\mu$, let
$q(x)=\sum_{k\in F}c_ke^{ik\cdot x}$, and let $\Phi\in\mathcal S(\R^d)$.
Then
\begin{equation}\label{eq:packet}
 \delta^d\big\|T_b(\Phi(\delta\,\cdot)q)
 -\Phi(\delta\,\cdot)T_{b|F}q\big\|_1\longrightarrow0
 \qquad(\delta\downarrow0).
\end{equation}
\end{lemma}

\begin{proof}
For one frequency, convolution by $\mu$ and the change of variables
$z=\delta x$ bound the normalized error by
\[
 \int_{\R^d}\|\Phi(\,\cdot-\delta y)-\Phi\|_1\,d|\mu|(y).
\]
This tends to zero by $L^1$ translation continuity and dominated
convergence, since the integrand is at most $2\|\Phi\|_1$.  Sum over the
finite set $F$.
\end{proof}

We also use the elementary periodic averaging formula.  If $a$ is
$2\pi$-periodic and integrable on $\T^d$, while $\Phi\in\mathcal S$ is
nonnegative, then
\begin{equation}\label{eq:averaging}
 \delta^d\int_{\R^d}\Phi(\delta x)|a(x)|\,dx
 \longrightarrow A_\Phi\int_{\T^d}|a(y)|\,dy,
 \qquad
 A_\Phi=(2\pi)^{-d}\int_{\R^d}\Phi.
\end{equation}
One proves this first for continuous $a$ by dividing $\R^d$ into period
cubes and taking a Riemann sum, then obtains the stated form by $L^1$
approximation.

\section{Restriction to the torus}

It suffices first to sample on $\Z^d$.  Let $P$ be a trigonometric
polynomial with zero constant term and finite spectrum
$F\subset\Z^d\setminus\{0\}$.  Choose a nonzero Schwartz function $\psi$
whose Fourier transform has compact support and put $\Phi=|\psi|^2$.  Thus
$\Phi\ge0$, $\Phi\in\mathcal S$, and $\wh\Phi$ is compactly supported.
Set
\[
 f_\delta(x)=\Phi(\delta x)P(x).
\]
For all sufficiently small $\delta$, the Fourier support of $f_\delta$ is
contained in a fixed compact subset of $\R^d\setminus\{0\}$.  Choose
$\chi\in C_c^\infty(\R^d\setminus\{0\})$ equal to one on that set and put
$b=m\chi$.  Lemma~\ref{lem:local-L1} applies, and $T_mf_\delta=T_bf_\delta$.

The product rule, \eqref{eq:averaging}, and the extra factor $\delta$ in the
derivative of the envelope give
\begin{equation}\label{eq:input-limit}
 \delta^d\|f_\delta\|_{\dot W_{\R}}
 \longrightarrow A_\Phi\|P\|_{\dot W_{\T}}.
\end{equation}
Indeed,
\[
 \partial_jf_\delta
 =\Phi(\delta x)\partial_jP(x)
   +\delta(\partial_j\Phi)(\delta x)P(x),
\]
and the normalized $L^1$ norm of the second term is $O(\delta)$.

Let $\gamma_k=m(k)$.  Commutation with derivatives gives
\[
 \partial_jT_mf_\delta
 =T_b\big(\Phi(\delta\,\cdot)\partial_jP\big)
  +\delta T_b\big((\partial_j\Phi)(\delta\,\cdot)P\big).
\]
The last term again has normalized norm $O(\delta)$ because convolution by
the finite measure representing $b$ is $L^1$ bounded.  Lemma~\ref{lem:packet}
applied to $\partial_jP$, followed by \eqref{eq:averaging}, yields
\begin{equation}\label{eq:output-limit}
 \delta^d\|T_mf_\delta\|_{\dot W_{\R}}
 \longrightarrow A_\Phi\|T_\gamma P\|_{\dot W_{\T}}.
\end{equation}
Combining \eqref{eq:input-limit}, \eqref{eq:output-limit}, and the Euclidean
multiplier inequality proves
\[
 \|T_\gamma P\|_{\dot W_{\T}}
 \le\|m\|_{\M(\dot W_{\R})}\|P\|_{\dot W_{\T}}.
\]
Trigonometric polynomials modulo constants are dense, so $\gamma$ extends to
the asserted torus multiplier.  Finally, the symbol $m_\eps(\xi)=m(\eps\xi)$
has the same Euclidean multiplier norm as $m$ by dilation.  Applying the
unit-lattice result to $m_\eps$ proves part~1 of
Theorem~\ref{thm:main}.

\section{Approximation from shrinking lattices}

Assume $\eps_n\downarrow0$ and
$\|\gamma^{\eps_n}\|_{\M(\dot W_{\T})}\le C$.  Testing on the nonconstant
exponentials shows that $|m(\eps_nk)|\le C$ for every $k\ne0$.  Given
$\xi\in\R^d$, choose nonzero $k_n\in\Z^d$ with $\eps_nk_n\to\xi$; continuity
then gives $|m(\xi)|\le C$.  Thus $m$ is bounded, as required in the
multiplier definition.  We first take
$f\in\mathcal S(\R^d)$ with $\wh f\in C_c^\infty(\R^d)$.  Put
$L_n=2\pi/\eps_n$ and define the finite trigonometric polynomial
\begin{equation}\label{eq:Pn}
 P_n(y)=\sum_{k\in\Z^d}
 \left(\frac{\eps_n}{2\pi}\right)^d
 \wh f(\eps_nk)e^{ik\cdot y}.
\end{equation}
Poisson summation gives, on the central cube $Q_{L_n}$ of side $L_n$,
\begin{equation}\label{eq:periodization}
 P_n(\eps_nx)=F_n(x):=\sum_{\ell\in\Z^d}f(x+L_n\ell).
\end{equation}
For each derivative, the period cubes partition $\R^d$.  Hence the triangle
inequality gives the upper bound by the total Euclidean norm, while the
reverse triangle inequality on the central summand loses only the exterior
tail.  Therefore
\begin{equation}\label{eq:input-period}
 \sum_{j=1}^d\int_{Q_{L_n}}|\partial_jF_n(x)|\,dx
 \longrightarrow\|f\|_{\dot W_{\R}}.
\end{equation}

The rescaled output derivatives are
\begin{align}
 H_{n,j}(x)
 &:=\eps_n(\partial_{y_j}T_{\gamma^{\eps_n}}P_n)(\eps_nx)\notag\\
 &=\sum_{k\in\Z^d}\left(\frac{\eps_n}{2\pi}\right)^d
 i\eps_nk_jm(\eps_nk)\wh f(\eps_nk)e^{i\eps_nk\cdot x}.
 \label{eq:Riemann}
\end{align}
For every fixed compact set of $x$'s, this is a Riemann sum converging
uniformly to
\begin{equation}\label{eq:limit-gradient}
 G_j(x)=(2\pi)^{-d}\int_{\R^d}
 i\xi_jm(\xi)\wh f(\xi)e^{ix\cdot\xi}\,d\xi
 =\partial_jT_mf(x).
\end{equation}
Indeed, the integrand is continuous and compactly supported, and its modulus
of continuity in $\xi$ is uniform for $x$ in a fixed compact set.

Changing variables $y=\eps_nx$ in the torus multiplier inequality cancels
the common factor $\eps_n^{1-d}$ and gives
\begin{equation}\label{eq:scaled-bound}
 \sum_j\int_{Q_{L_n}}|H_{n,j}(x)|\,dx
 \le C\sum_j\int_{Q_{L_n}}|\partial_jF_n(x)|\,dx.
\end{equation}
Fix a ball $B_R$.  For large $n$, it lies in $Q_{L_n}$.  Uniform convergence
in \eqref{eq:Riemann}, \eqref{eq:input-period}, and
\eqref{eq:scaled-bound} show that
\[
 \sum_j\int_{B_R}|G_j(x)|\,dx\le C\|f\|_{\dot W_{\R}}.
\]
Letting $R\to\infty$ proves
\begin{equation}\label{eq:dense-bound}
 \|T_mf\|_{\dot W_{\R}}\le C\|f\|_{\dot W_{\R}}
\end{equation}
on the band-limited Schwartz class.

For completeness, that class is dense in $\dot W^{1,1}(\R^d)$ modulo
constants when $d\ge2$.  The endpoint Sobolev inequality supplies a
representative $u\in L^{d/(d-1)}$ for every homogeneous class.  If
$\eta_R(x)=\eta(x/R)$ is a standard cutoff, then
\[
 \|(1-\eta_R)\nabla u\|_1\to0,
 \qquad
 \|u\nabla\eta_R\|_1
 \lesssim\|u\|_{L^{d/(d-1)}(\{R\lesssim|x|\lesssim2R\})}\to0.
\]
Mollification therefore proves density of $C_c^\infty$.  Convolving a
compactly supported smooth function with an approximate identity
$\rho_t\in\mathcal S$ for which $\wh\rho_t$ is compactly supported gives
band-limited Schwartz approximants and convergence of every first derivative
in $L^1$.  Thus \eqref{eq:dense-bound} extends uniquely to all of
$\dot W_{\R}$, proving part~2.  Combining the two parts proves
\eqref{eq:norm-identity}.

\section{Scope and novelty note}

The theorem is complete for the endpoint first-order homogeneous space
studied in \cite{KW}: restriction, shrinking-lattice approximation, and
exact norm equality.  It does not settle all orders and exponents; the
low-degree polynomial quotients and test-class density across homogeneous
Sobolev thresholds require separate analysis.

A literature screen through 12 August 2026 found later continuity results on
$\dot W^{l,1}$ and $\dot W^{l,\infty}$ \cite{Curca}, and related transference
for the \emph{inhomogeneous} $W^{l,1}$ spaces \cite{CW}.  Neither source
states the homogeneous shrinking-lattice theorem above.  This is a bounded,
not exhaustive, novelty screen.

\small
\begin{thebibliography}{9}
\bibitem{KW}
K.~Kazaniecki and M.~Wojciechowski,
\emph{On the continuity of Fourier multipliers on the homogeneous Sobolev
spaces $\dot W^1_1(\R^d)$}, Ann. Inst. Fourier 66 (2016), no.~3,
1247--1260; arXiv:1306.1437.

\bibitem{Rudin}
W.~Rudin, \emph{Fourier Analysis on Groups}, Wiley, 1962.

\bibitem{Curca}
E.~Curca, \emph{On the continuity of Fourier multipliers on
$\dot W^{l,1}(\R^d)$ and $\dot W^{l,\infty}(\R^d)$}, J. Funct. Anal.
(2022), 109573.

\bibitem{CW}
E.~Curca and M.~Wojciechowski,
\emph{Nonexistence of Henkin type projections via a Wiener theorem for
multipliers}, arXiv:2604.23161 (2026).
\end{thebibliography}

\end{document}
